%% ======================================================================
\chapter{Further topics and recent advances}
\label{chap:advances}
%% ======================================================================

This last chapter is a reading guide. It briefly presents some active research
directions, with pointers to the literature; it is a good starting point for
the ``research'' part of the project.

\section{Beyond bounded losses and i.i.d.\ data}

All the bounds of this course assume a loss in $[0,1]$ and i.i.d.\ data. For
unbounded losses (squared loss, log-loss), the term $\mathcal I_\Delta(m)$ of
Theorem~\ref{thm:general} must be controlled through assumptions on the tails of the
loss (sub-Gaussian, sub-exponential, bounded moments):
see \citet{alquier2016properties}, \citet{germain2016pac} and the recent
unified treatments of \citet{rodriguez2024more} and \citet{chugg2023unified}, which
also give \emph{anytime-valid} bounds (holding simultaneously for all sample
sizes, a useful property for online learning and sequential decisions).
Martingale techniques extend PAC-Bayes to dependent data and to reinforcement
learning \citep{seldin2012pac,chugg2023unified}.

\section{Other divergences and priors}

The KL divergence comes from the change of measure \eqref{eq:dv}; other
changes of measure yield bounds with other divergences, e.g.\ Rényi
divergences \citep{begin2016pac}, which can be much smaller when $Q$ and $P$
have different supports. Tighter bounds are also obtained with Bernstein-type
inequalities \citep{tolstikhin2013pac,mhammedi2019pac} and, for the
ternary-valued losses that appear in majority votes, with the split-kl
inequality \citep{wu2022split}. On the prior side, data-dependent priors
remain one of the most effective levers
\citep{parrado2012pac,dziugaite2018data,perez2021tighter}.

\section{Information-theoretic generalization bounds}

PAC-Bayesian bounds control the average risk under a posterior that depends on
$S$. A close family of results bounds the \emph{expected} generalization gap of
a (randomized) algorithm by the \emph{mutual information} $I(S;h)$ between the
sample and the output \citep{xu2017information}. Averaging a PAC-Bayesian bound
over $S$ with the ``optimal'' prior $P=\E_SQ_S$ turns $\E_S\KL(Q_S\|P)$ into
$I(S;h)$, which shows the deep connection between the two approaches. The
monograph by \citet{hellstrom2025generalization} presents both points of view in a
unified way.

\section{Deterministic predictors: disintegration}

A drawback of PAC-Bayes is that it certifies a \emph{randomized} predictor.
Disintegrated bounds \citep{blanchard2007occam,rivasplata2020pac} hold with
high probability over the joint draw of $S$ and of a single $h\sim Q_S$, so
that the one model actually deployed is certified;
\citet{viallard2024general} give a general framework and practical algorithms
for such bounds. For majority votes, this line of work leads to bounds on the
vote itself rather than on the Gibbs classifier.

\section{Domain adaptation and transfer}

When the training distribution (source) differs from the test distribution
(target), the risk of a majority vote on the target can be related to its risk
on the source and to a \emph{divergence between domains}. PAC-Bayes is
particularly well suited because the disagreement $\dD(Q)$ can be computed on
\emph{unlabeled} target data: \citet{germain2013pac,germain2020pac} define a
domain divergence as the difference of disagreements between source and target
and derive PAC-Bayesian domain adaptation bounds and algorithms. Similar ideas
are used for multi-view learning and for fairness constraints.

\section{Open questions}

\begin{itemize}
\item \textbf{Tightness vs.\ learning.} The tightest certificates (kl, CCPBB) are
      not always the best learning objectives; understanding which bound to
      minimize to obtain the best \emph{predictor} is still open.
\item \textbf{Large ensembles and large models.} How to design priors for boosted
      ensembles with thousands of trees, or for ensembles of large neural
      networks, so that bounds remain non-vacuous \citep{lotfi2022pac}?
\item \textbf{Diversity.} Second-order bounds formalize the role of diversity;
      can we design ensemble algorithms that explicitly optimize the
      strength--diversity trade-off with guarantees, beyond random forests?
\end{itemize}

\section*{Further reading}
\addcontentsline{toc}{section}{Further reading}
\begin{itemize}
\item Introductions: \citet{langford2005tutorial}, \citet{guedj2019primer},
      \citet{alquier2024user}.
\item Majority votes: \citet{germain2015risk}, \citet{masegosa2020second},
      \citet{wu2021chebyshev}, \citet{viallard2021self}.
\item Advanced: \citet{catoni2007pac}, \citet{hellstrom2025generalization}.
\end{itemize}
